% ------------------------------------------------------------------------------
Here's a basic job script, \file{tcsh.sh} shown in \xf{fig:tcsh.sh}.
You can copy it from our \href{https://github.com/NAG-DevOps/speed-hpc}{GitHub repository}.

\begin{figure}[htpb]
	\lstinputlisting[language=csh,frame=single,basicstyle=\ttfamily]{tcsh.sh}
	\caption{Source code for \file{tcsh.sh}}
	\label{fig:tcsh.sh}
\end{figure}

\noindent The first line is the shell declaration (also know as a shebang) and sets the shell to \emph{tcsh}.
The lines that begin with \texttt{\#SBATCH} are directives for the scheduler.
\begin{itemize}
	%\item \texttt{-N} sets \emph{qsub-test} as the jobname
	\item \texttt{-J} (or \option{--job-name}) sets \emph{tcsh-test} as the job name.
	%\item \texttt{-cwd} tells the scheduler to execute the job from the current working directory
	%\item \texttt{--chdir} tells the scheduler to execute the job from the current working directory
	%\item \texttt{-l h\_vmem=1GB} requests and assigns 1GB of memory to the job. CPU jobs \emph{require} the \texttt{-l h\_vmem} option to be set.
	\item \texttt{--mem=1GB} requests and assigns 1GB of memory to the job. 
	Jobs require the \texttt{--mem} option to be set either in the script
	or on the command line; \textbf{if it's missing, job submission will be rejected.}
\end{itemize}

\noindent The script then:
\begin{enumerate}
	\item Sleeps on a node for 30 seconds
	\item Uses the \tool{module} command to load the \texttt{gurobi/8.1.0} environment
	\item Prints the list of loaded modules into a file
\end{enumerate}

%The scheduler command, \tool{qsub}, is used to submit (non-interactive) jobs. 
\noindent The scheduler command, \tool{sbatch}, is used to submit (non-interactive) jobs. 
%From an ssh session on speed-submit, submit this job with \texttt{qsub ./tcsh.sh}.
From an ssh session on \texttt{speed-submit}, submit this job with \texttt{sbatch ./tcsh.sh}.
%You will see, \texttt{"Your job X ("qsub-test") has been submitted"}.
You will see, \texttt{"Submitted batch job 2653"} where $2653$ is a job ID assigned.
%The command, \tool{qstat}, can be used 
The commands, \tool{squeue} and \tool{sinfo} can be used 
%to look at the status of the cluster: \texttt{qstat -f -u "*"}.
to look at the status of the cluster: \texttt{squeue}.
You will see something like this: 

%\small
%\begin{verbatim}
%queuename                      qtype resv/used/tot. load_avg arch          states
%---------------------------------------------------------------------------------
%a.q@speed-01.encs.concordia.ca BIP   0/0/32         0.01     lx-amd64
%---------------------------------------------------------------------------------
%a.q@speed-03.encs.concordia.ca BIP   0/0/32         0.01     lx-amd64
%---------------------------------------------------------------------------------
%a.q@speed-25.encs.concordia.ca BIP   0/0/32         0.01     lx-amd64
%---------------------------------------------------------------------------------
%a.q@speed-27.encs.concordia.ca BIP   0/0/32         0.01     lx-amd64
%---------------------------------------------------------------------------------
%g.q@speed-05.encs.concordia.ca BIP   0/0/32         0.02     lx-amd64
     %144   100.00000 qsub-test nul-uge     r     12/03/2018 16:39:30    1 
     %62624 0.09843 case_talle x_yzabc      r     11/09/2021 16:50:09    32
%---------------------------------------------------------------------------------
%g.q@speed-17.encs.concordia.ca BIP   0/0/32         0.01     lx-amd64
%---------------------------------------------------------------------------------
%s.q@speed-07.encs.concordia.ca BIP   0/0/32         0.04     lx-amd64
%---------------------------------------------------------------------------------
%s.q@speed-08.encs.concordia.ca BIP   0/0/32         0.01     lx-amd64
%---------------------------------------------------------------------------------
%s.q@speed-09.encs.concordia.ca BIP   0/0/32         0.01     lx-amd64
%---------------------------------------------------------------------------------
%s.q@speed-10.encs.concordia.ca BIP   0/32/32        32.72    lx-amd64
     %62624 0.09843 case_talle x_yzabc      r     11/09/2021 16:50:09    32
%---------------------------------------------------------------------------------
%s.q@speed-11.encs.concordia.ca BIP   0/32/32        32.08    lx-amd64
     %62679 0.14212 CWLR_DF    a_bcdef      r     11/10/2021 17:25:19    32
%---------------------------------------------------------------------------------
%s.q@speed-12.encs.concordia.ca BIP   0/32/32        32.10    lx-amd64
     %62749 0.09000 CLOUDY     z_abc        r     11/11/2021 21:58:12    32
%---------------------------------------------------------------------------------
%s.q@speed-15.encs.concordia.ca BIP   0/4/32         0.03     lx-amd64
     %62753 82.47478 matlabLDPa b_bpxez      r     11/12/2021 08:49:52     4
%---------------------------------------------------------------------------------
%s.q@speed-16.encs.concordia.ca BIP   0/32/32        32.31    lx-amd64
     %62751 0.09000 CLOUDY     z_abc        r     11/12/2021 06:03:54    32
%---------------------------------------------------------------------------------
%s.q@speed-19.encs.concordia.ca BIP   0/32/32        32.22    lx-amd64
%---------------------------------------------------------------------------------
%...
%---------------------------------------------------------------------------------
%s.q@speed-35.encs.concordia.ca BIP   0/32/32        2.78     lx-amd64
     %62754 7.22952 qlogin-tes a_tiyuu      r     11/12/2021 10:31:06    32
%---------------------------------------------------------------------------------
%s.q@speed-36.encs.concordia.ca BIP   0/0/32         0.03     lx-amd64
%etc.
\small
\begin{verbatim}
[serguei@speed-submit src] % squeue -l
Thu Oct 19 11:38:54 2023
JOBID PARTITION     NAME     USER    STATE       TIME TIME_LIMI  NODES NODELIST(REASON)
 2641        ps interact   b_user  RUNNING   19:16:09 1-00:00:00      1 speed-07
 2652        ps interact   a_user  RUNNING      41:40 1-00:00:00      1 speed-07
 2654        ps tcsh-tes  serguei  RUNNING       0:01 7-00:00:00      1 speed-07
[serguei@speed-submit src] % sinfo
PARTITION AVAIL  TIMELIMIT  NODES  STATE NODELIST
ps*          up 7-00:00:00     14  drain speed-[08-10,12,15-16,20-22,30-32,35-36]
ps*          up 7-00:00:00      1    mix speed-07
ps*          up 7-00:00:00      7   idle speed-[11,19,23-24,29,33-34]
pg           up 1-00:00:00      1  drain speed-17
pg           up 1-00:00:00      3   idle speed-[05,25,27]
pt           up 7-00:00:00      7   idle speed-[37-43]
pa           up 7-00:00:00      4   idle speed-[01,03,25,27]
\end{verbatim}
\normalsize

Remember that you only have 30 seconds before the job is essentially over, so 
if you do not see a similar output, either adjust the sleep time in the 
%script, or execute the \tool{qstat} statement more quickly. The \tool{qstat} 
script, or execute the \tool{sbatch} statement more quickly. The \tool{squeue} 
output listed above shows you that your job is running on node \texttt{speed-07}, 
that it has a job number of 2654, its time limit of 7 days, etc.
% TODO
%, that it 
%was started at 16:39:30 on 12/03/2018, and that it is a single-core job (the 
%default). 

Once the job finishes, there will be a new file in the directory that the job 
%was started from, with the syntax of, \texttt{"job name".o"job number"}, so 
was started from, with the syntax of, \texttt{slurm-"job id".out}, so 
%in this example the file is, qsub \file{test.o144}. This file represents the 
in this example the file is, \file{slurm-2654.out}. This file represents the 
standard output (and error, if there is any) of the job in question. If you 
look at the contents of your newly created file, you will see that it 
contains the output of the, \texttt{module list} command. 
Important information is often written to this file.
%
%Congratulations on your first job!